\chapter{The Gaussian Chain by Linear Algebra}\label{ch:gaussian}

The stationary Gaussian AR(1) chain is the one model in this thesis whose joint
score can be written down in closed form and checked by hand. This chapter does
that with matrices: the clean law, its dense covariance and sparse precision,
the score in two equivalent forms, and the lifecycle the coupling structure
passes through as the diffusion time grows. 

\section{The joint law and the object of study}\label{sec:g-model}

A data point is a trajectory $a = (a_0, \dots, a_{L-1})$ of the stationary
AR(1) chain \eqref{eq:ar1} with $\sigma_\eta^2 = 1 - \alpha^2$. Each frame
passes independently through the OU channel \eqref{eq:ou-channel} for a
time $t$,
\begin{equation}
  x_k \;=\; e^{-t} a_k + \sqrt{\Deltat}\;\xi_k,
  \qquad \xi_k \sim \Nd(0,1)\ \text{i.i.d.},
  \label{eq:g-channel}
\end{equation}
and the object of study is the joint score of the noisy sequence,
\begin{equation}
  \score(x, t) \;=\; \nabla_x \log P_t(x) \;\in\; \R^L,
  \qquad
  P_t(x) = \int_{\R^L} P_0(a) \prod_{k=0}^{L-1}
    \Nd\big(x_k;\, e^{-t}a_k,\, \Deltat\big)\, \dd a .
  \label{eq:g-object}
\end{equation}
The gradient is taken with respect to all $L$ noisy frames simultaneously,
under the law of the whole noisy trajectory
(Section~\ref{sec:tweedie}).

\section{The clean law: from the recursion to the covariance}
\label{sec:g-clean}

\subsection{Unrolling the recursion}

The AR(1) chain \eqref{eq:ar1} is the first-order recursion
\begin{equation}
  a_{k+1}
  = \alpha a_k + \eta_k,
  \qquad
  \eta_k \stackrel{\mathrm{iid}}{\sim} \Nd(0,\sigma_\eta^2),
  \qquad
  \sigma_\eta^2 = 1-\alpha^2,
  \label{eq:g-recursion}
\end{equation}
with
\[
  a_0\sim\Nd(0,1),
  \qquad
  a_0 \perp (\eta_0,\eta_1,\ldots),
  \qquad
  |\alpha|<1.
\]
The parameter $\alpha$ controls how strongly one frame is inherited by the
next: $\alpha=0$ gives independent frames, while $|\alpha|$ close to one
produces long persistence along the trajectory.

The recursion can be expanded explicitly. The first step gives
\[
  a_1 = \alpha a_0+\eta_0.
\]
Substituting this expression into the next step,
\[
  a_2
  =\alpha a_1+\eta_1
  =\alpha(\alpha a_0+\eta_0)+\eta_1
  =\alpha^2a_0+\alpha\eta_0+\eta_1.
\]
Repeating once more,
\[
  a_3
  =\alpha a_2+\eta_2
  =\alpha^3a_0+\alpha^2\eta_0+\alpha\eta_1+\eta_2.
\]
The pattern is therefore
\begin{equation}
  \boxed{
  a_k
  =\alpha^k a_0
  +\sum_{j=0}^{k-1}\alpha^{\,k-1-j}\eta_j
  }
  \qquad (k\geq1).
  \label{eq:g-unroll}
\end{equation}
Each innovation is propagated forward by one factor of $\alpha$ for every
subsequent step: the most recent innovation $\eta_{k-1}$ enters with
coefficient one, whereas an innovation generated $m$ steps earlier enters
with coefficient $\alpha^m$.

Equation~\eqref{eq:g-unroll} follows formally by induction. Suppose it holds
at step $k$. Then
\begin{align}
  a_{k+1}
  &=\alpha a_k+\eta_k \\
  &=\alpha
    \left(
      \alpha^k a_0
      +\sum_{j=0}^{k-1}\alpha^{\,k-1-j}\eta_j
    \right)
    +\eta_k \\
  &=\alpha^{k+1}a_0
    +\sum_{j=0}^{k-1}\alpha^{\,k-j}\eta_j
    +\eta_k \\
  &=\alpha^{k+1}a_0
    +\sum_{j=0}^{k}\alpha^{\,k-j}\eta_j,
\end{align}
which is exactly \eqref{eq:g-unroll} with $k+1$ in place of $k$.

This representation already exposes the probabilistic structure of the
whole trajectory. Define
\[
  \varepsilon
  =
  (a_0,\eta_0,\ldots,\eta_{L-2})^{\mathsf T}.
\]
Its entries are independent Gaussian variables, so $\varepsilon$ is jointly
Gaussian. Equation~\eqref{eq:g-unroll} shows that there exists a deterministic
lower-triangular matrix $G_\alpha$ such that
\begin{equation}
  a
  =
  (a_0,\ldots,a_{L-1})^{\mathsf T}
  =
  G_\alpha\,\varepsilon.
  \label{eq:g-linear-map}
\end{equation}
A deterministic linear transformation of a Gaussian vector is Gaussian;
hence the complete trajectory $a$ is jointly Gaussian.

The choice $\sigma_\eta^2=1-\alpha^2$ makes this Gaussian chain stationary
with unit marginal variance. Indeed, if $\operatorname{Var}(a_k)=1$, independence of
$a_k$ and $\eta_k$ gives
\begin{align}
  \operatorname{Var}(a_{k+1})
  &=\operatorname{Var}(\alpha a_k+\eta_k) \\
  &=\alpha^2\operatorname{Var}(a_k)+\operatorname{Var}(\eta_k) \\
  &=\alpha^2+(1-\alpha^2)
  =1.
\end{align}
Since $\operatorname{Var}(a_0)=1$, induction gives
\begin{equation}
  a_k\sim\Nd(0,1)
  \qquad\text{for every }k.
  \label{eq:g-stationary-marginal}
\end{equation}

Finally, \eqref{eq:g-unroll} makes the loss of memory explicit. Conditional
on the initial state,
\begin{equation}
  \E[a_k\mid a_0]
  =\alpha^k a_0,
  \label{eq:g-memory}
\end{equation}
because every innovation has zero mean. Since $|\alpha|<1$,
\[
  |\alpha|^k\longrightarrow0
  \qquad\text{as }k\longrightarrow\infty.
\]
Thus the influence of the initial condition decays geometrically. More
precisely, for $j\leq k$, iterating the recursion only from frame $j$ gives
\[
  a_k
  =\alpha^{k-j}a_j
   +\sum_{r=j}^{k-1}\alpha^{\,k-1-r}\eta_r.
\]
The second term is independent of $a_j$, and therefore
\begin{equation}
  \operatorname{Cov}(a_j,a_k)
  =\alpha^{k-j}\operatorname{Var}(a_j)
  =\alpha^{k-j}.
\end{equation}
By symmetry,
\begin{equation}
  \boxed{
  \operatorname{Cov}(a_j,a_k)=\alpha^{|k-j|}
  }.
  \label{eq:g-ar1-covariance}
\end{equation}
The same geometric factor therefore controls both notions of memory:
the effect of the initial condition and the correlation between frames
separated by a given lag.


\subsection{The stationary covariance is Toeplitz}

With mean zero, the joint law is fixed by the covariances
$\operatorname{Cov}(a_i,a_j)$. Take $j = k$ first. The terms of
\eqref{eq:g-unroll} are independent, so their variances add:
\begin{equation}
  \operatorname{Var}(a_k)
  = \alpha^{2k}\underbrace{\operatorname{Var}(a_0)}_{1}
  + \sigma_\eta^2 \sum_{j=0}^{k-1}\alpha^{\,2(k-1-j)}
  = \alpha^{2k}
  + \sigma_\eta^2 \sum_{i=0}^{k-1}\big(\alpha^{2}\big)^{i},
  \label{eq:g-var-sum}
\end{equation}
the last step substituting $i = k-1-j$. The remaining sum is geometric with
ratio $\alpha^2 \neq 1$,
\begin{equation}
  \sum_{i=0}^{k-1}\big(\alpha^{2}\big)^{i}
  = \frac{1 - \alpha^{2k}}{1 - \alpha^{2}},
  \label{eq:g-geom}
\end{equation}
and with $\sigma_\eta^2 = 1-\alpha^2$ the prefactor cancels the denominator:
\begin{equation}
  \operatorname{Var}(a_k)
  = \alpha^{2k} + (1-\alpha^2)\,\frac{1-\alpha^{2k}}{1-\alpha^2}
  = \alpha^{2k} + \big(1 - \alpha^{2k}\big)
  = 1 ,
  \label{eq:g-var-one}
\end{equation}
for every $k$: the pairing $a_0 \sim \Nd(0,1)$, $\sigma_\eta^2 = 1-\alpha^2$ is
exactly what holds the marginal variance constant along the chain. For the
cross-covariance, split \eqref{eq:g-unroll} at index $k$,
\begin{equation}
  a_{k+d}
  = \alpha^{d} a_k
  + \sum_{j=k}^{k+d-1}\alpha^{\,k+d-1-j}\eta_j
  \qquad (d \ge 0),
  \label{eq:g-split}
\end{equation}
where the sum uses only $\eta_k,\dots,\eta_{k+d-1}$, all independent of $a_k$.
Hence
\begin{equation}
  \operatorname{Cov}(a_k, a_{k+d})
  = \alpha^{d}\,\operatorname{Var}(a_k)
  = \alpha^{d},
  \qquad\text{i.e.}\qquad
  (\Sigz)_{ij} = \alpha^{|i-j|}.
  \label{eq:g-cov-lag}
\end{equation}

Collecting mean and covariance, the clean joint density is
\begin{equation}
  P_0(a)
  = \Nd(a_0; 0, 1)\prod_{k=0}^{L-2}
      \Nd\big(a_{k+1};\, \alpha a_k,\, \sigma_\eta^2\big)
  = \Nd(a;\, 0,\, \Sigz),
  \label{eq:g-P0}
\end{equation}
the product being the chain factorisation \eqref{eq:chain-factorisation} and
the right-hand side its covariance form, with
\begin{equation}
  \Sigz
  =
  \begin{pmatrix}
    1 & \alpha & \alpha^2 & \cdots & \alpha^{L-1} \\
    \alpha & 1 & \alpha & \cdots & \alpha^{L-2} \\
    \alpha^2 & \alpha & 1 & \cdots & \alpha^{L-3} \\
    \vdots & \vdots & \vdots & \ddots & \vdots \\
    \alpha^{L-1} & \alpha^{L-2} & \alpha^{L-3} & \cdots & 1
  \end{pmatrix}.
  \label{eq:g-Sig0-matrix}
\end{equation}

\begin{proposition}[Stationary covariance]\label{prop:g-Sigma0}
$(\Sigz)_{ij} = \alpha^{|i-j|}$: a symmetric Toeplitz matrix ---
constant along every diagonal --- with geometrically decaying bands,
autocovariance $\gamma(d) = \alpha^{|d|}$. It is dense for $\alpha \neq 0$;
at $\alpha = 0$ the chain is white and $\Sigz = I$.
\end{proposition}

The Toeplitz form is the matrix print of weak stationarity: by
\eqref{eq:weak-stationary}--\eqref{eq:toeplitz} the covariance of a stationary
sequence depends only on the index difference, so the $d$-th diagonal of
\eqref{eq:g-Sig0-matrix} carries the single number $\gamma(d) = \alpha^d$.
Three readings of that line are used later, all for the coupled case
$\alpha \neq 0$ (at $\alpha = 0$ the chain is white and $\Sigz = I$). The
matrix is \emph{dense}: every pair of frames is correlated and no entry of
$\Sigz$ is exactly zero. The correlation magnitude \emph{decays
geometrically}, $|\gamma(d)| = e^{-d/\tau}$ with correlation length
$\tau = -1/\ln|\alpha|$ --- for $\alpha < 0$ the sign alternates with the
lag while the modulus follows the same law --- so far-apart frames are weakly
but not trivially coupled. And the coupling is \emph{global though the
mechanism is local}: \eqref{eq:g-recursion} touches only neighbours, yet its
covariance reaches across the whole chain. The locality is still there --- in
the inverse.

\section{The clean law: the precision is tridiagonal}
\label{sec:g-precision}

\subsection{Reading the precision off a short chain}

The score will depend not on $\Sigz$ but on its inverse, the \emph{precision}
$\Qz := \Sigz^{-1}$. For short chains it is an exercise. At $L = 2$,
\begin{equation}
  \Sigz = \begin{pmatrix} 1 & \alpha \\ \alpha & 1 \end{pmatrix},
  \qquad
  \Qz
  = \frac{1}{1-\alpha^2}\begin{pmatrix} 1 & -\alpha \\ -\alpha & 1 \end{pmatrix}
  = \frac{1}{\sigma_\eta^2}\begin{pmatrix} 1 & -\alpha \\ -\alpha & 1 \end{pmatrix},
  \label{eq:g-Q0-K2}
\end{equation}
from the $2\times2$ rule
$\begin{psmallmatrix}a&b\\c&d\end{psmallmatrix}^{-1}
 = (ad-bc)^{-1}\begin{psmallmatrix}d&-b\\-c&a\end{psmallmatrix}$
with determinant $1-\alpha^2$. At $L = 3$ the cofactor inverse gives
\begin{equation}
  \Sigz = \begin{pmatrix}
    1 & \alpha & \alpha^2 \\
    \alpha & 1 & \alpha \\
    \alpha^2 & \alpha & 1
  \end{pmatrix},
  \qquad
  \Qz = \frac{1}{\sigma_\eta^2}\begin{pmatrix}
    1 & -\alpha & 0 \\
    -\alpha & 1+\alpha^2 & -\alpha \\
    0 & -\alpha & 1
  \end{pmatrix}
  \label{eq:g-Q0-K3}
\end{equation}
(one line of $\Qz\Sigz$ confirms it). Two features of \eqref{eq:g-Q0-K3} hold
for every $L$: the corner entry $(\Qz)_{0,2}$ is \emph{exactly} zero, and the
interior diagonal carries an extra $\alpha^2$ that the two endpoints do not.

\subsection{The tridiagonal precision and what its zeros mean}

\begin{proposition}[Tridiagonal precision]\label{prop:g-Q0}
The precision matrix $\Qz=\Sigz^{-1}$ is
\begin{equation}
  \boxed{
  \Qz
  =
  \frac{1}{\sigma_\eta^2}
  \begin{pmatrix}
    1          & -\alpha       &              &              & 0\\
    -\alpha    & 1+\alpha^2    & -\alpha      &              & \\
               & \ddots        & \ddots       & \ddots       & \\
               &               & -\alpha      & 1+\alpha^2   & -\alpha\\
    0          &               &              & -\alpha      & 1
  \end{pmatrix}.
  }
  \label{eq:g-Q0}
\end{equation}
Equivalently,
\begin{equation}
  (\Qz)_{kk}
  =
  \begin{cases}
    \dfrac{1}{\sigma_\eta^2},
      & k\in\{0,L-1\},\\[7pt]
    \dfrac{1+\alpha^2}{\sigma_\eta^2},
      & 0<k<L-1,
  \end{cases}
  \qquad
  (\Qz)_{k,k+1}
  =
  -\frac{\alpha}{\sigma_\eta^2},
\end{equation}
with $(\Qz)_{ij}=0$ whenever $|i-j|\geq2$.
\end{proposition}

There is no need to invert $\Sigz$ entry by entry. For a centred Gaussian
$a\sim\Nd(0,\Sigz)$,
\begin{equation}
  P_0(a)
  =
  (2\pi)^{-L/2}|\Sigz|^{-1/2}
  \exp\!\left(-\frac12 a^{\T}\Qz a\right),
  \qquad
  \Qz=\Sigz^{-1},
  \label{eq:g-gaussian-precision}
\end{equation}
and therefore
\begin{align}
  -2\log P_0(a)
  &=
  a^{\T}\Qz a
  +L\log(2\pi)+\log|\Sigz| \\
  &=
  a^{\T}\Qz a+\mathrm{const}.
  \label{eq:g-read-precision}
\end{align}
The precision matrix can thus be read directly from the quadratic part of
the factorised negative log-density.

\begin{toolbox}{Reading $\Qz$ off $-2\log P_0$}

The Markov factorisation \eqref{eq:g-P0} is
\begin{equation*}
  P_0(a)
  =
  p(a_0)
  \prod_{k=0}^{L-2}p(a_{k+1}\mid a_k),
\end{equation*}
with
\[
  a_0\sim\Nd(0,1),
  \qquad
  a_{k+1}\mid a_k
  \sim\Nd(\alpha a_k,\sigma_\eta^2).
\]
Discarding terms independent of $a$ gives
\begin{align}
  -2\log P_0(a)
  &=
  a_0^2
  +
  \frac{1}{\sigma_\eta^2}
  \sum_{k=0}^{L-2}
  \bigl(a_{k+1}-\alpha a_k\bigr)^2
  +\mathrm{const}
  \label{eq:g-precision-expand-1}\\
  &=
  a_0^2
  +
  \frac{1}{\sigma_\eta^2}
  \sum_{k=0}^{L-2}
  \left(
    a_{k+1}^2
    -2\alpha a_ka_{k+1}
    +\alpha^2a_k^2
  \right)
  +\mathrm{const}.
  \label{eq:g-precision-expand-2}
\end{align}
We compare this with the quadratic form of a symmetric matrix,
\begin{align}
  a^{\T}\Qz a
  &=
  \sum_{m=0}^{L-1}(\Qz)_{mm}a_m^2
  +
  2\sum_{0\leq m<n\leq L-1}
  (\Qz)_{mn}a_ma_n.
  \label{eq:g-general-quadratic-form}
\end{align}

\medskip
\noindent\textbf{Off-diagonal entries.}
The only products of distinct coordinates are
\begin{align}
  -\frac{2\alpha}{\sigma_\eta^2}
  \sum_{k=0}^{L-2}a_ka_{k+1}.
\end{align}
Matching coefficients with
$2(\Qz)_{k,k+1}a_ka_{k+1}$ gives
\begin{align}
  (\Qz)_{k,k+1}
  &=
  (\Qz)_{k+1,k}
  =
  -\frac{\alpha}{\sigma_\eta^2},
  \qquad k=0,\ldots,L-2,
  \label{eq:g-precision-offdiag}\\
  (\Qz)_{ij}
  &=0,
  \qquad |i-j|\geq2.
  \label{eq:g-precision-distant-zero}
\end{align}
The precision is therefore tridiagonal.

\medskip
\noindent\textbf{Interior diagonal.}
For $1\leq m\leq L-2$, the term $a_m^2$ appears in two adjacent transition
factors:
\begin{align}
  \frac{1}{\sigma_\eta^2}\bigl(a_m-\alpha a_{m-1}\bigr)^2
  &= \frac{1}{\sigma_\eta^2}a_m^2 - \frac{2\alpha}{\sigma_\eta^2}a_m a_{m-1} + \frac{\alpha^2}{\sigma_\eta^2}a_{m-1}^2,\\
  \frac{1}{\sigma_\eta^2}\bigl(a_{m+1}-\alpha a_m\bigr)^2
  &= \frac{\alpha^2}{\sigma_\eta^2}a_m^2 - \frac{2\alpha}{\sigma_\eta^2}a_{m+1}a_m + \frac{1}{\sigma_\eta^2}a_{m+1}^2.
\end{align}
Summing the coefficients of $a_m^2$ yields
\begin{equation}
  (\Qz)_{mm}
  =
  \frac{1+\alpha^2}{\sigma_\eta^2},
  \qquad
  1\leq m\leq L-2.
  \label{eq:g-precision-interior}
\end{equation}

\medskip
\noindent\textbf{Boundaries.}
The left endpoint also appears in its prior:
\begin{align}
  a_0^2
  +
  \frac{\alpha^2}{\sigma_\eta^2}a_0^2
  &=
  \frac{\sigma_\eta^2+\alpha^2}
       {\sigma_\eta^2}a_0^2\\
  &=
  \frac{1}{\sigma_\eta^2}a_0^2,
\end{align}
where the last equality uses
$\sigma_\eta^2=1-\alpha^2$. Thus
\begin{equation}
  (\Qz)_{00}
  =
  \frac{1}{\sigma_\eta^2}.
\end{equation}
The final coordinate appears only in the last transition,
\begin{align}
  \frac{1}{\sigma_\eta^2}
  \bigl(a_{L-1}-\alpha a_{L-2}\bigr)^2,
\end{align}
so similarly
\begin{equation}
  (\Qz)_{L-1,L-1}
  =
  \frac{1}{\sigma_\eta^2}.
\end{equation}

Putting the diagonal and off-diagonal coefficients together gives
\eqref{eq:g-Q0}. The calculation also explains the pattern: an interior
variable belongs to two transition factors, whereas an endpoint belongs to
only one; and no factor ever contains two non-neighbouring variables.
\hfill$\square$
\end{toolbox}

The zeros of \eqref{eq:g-Q0} have a direct probabilistic meaning. For a
jointly Gaussian vector,
\begin{equation}
  (\Qz)_{ij}=0
  \quad\Longleftrightarrow\quad
  a_i\perp a_j
  \mid
  a_{\setminus\{i,j\}}.
  \label{eq:g-precision-ci}
\end{equation}
Thus tridiagonality is the matrix expression of the first-order Markov
structure: once the neighbouring frames are known, an interior frame is
conditionally independent of all more distant frames.

The origin of the sparsity is more general than the Gaussian model itself.
For any smooth first-order Markov density,
\begin{align}
  -\log P(a)
  &=
  -\log p(a_0)
  -
  \sum_{k=0}^{L-2}
  \log p(a_{k+1}\mid a_k),
\end{align}
and each summand contains only neighbouring coordinates. Consequently,
\begin{equation}
  \frac{\partial^2[-\log P(a)]}
       {\partial a_i\,\partial a_j}
  =0,
  \qquad |i-j|\geq2.
  \label{eq:g-markov-hessian-sparsity}
\end{equation}
For a general Markov chain this is a sparse, possibly state-dependent
Hessian. In the Gaussian case the Hessian is constant and equals the
precision matrix $\Qz$.

The same matrix also tells us how one hidden frame is predicted from the
others. For a centred Gaussian,
\begin{equation}
  \E[a_k\mid a_{\setminus k}]
  =
  -(\Qz)_{kk}^{-1}
  \sum_{l\neq k}(\Qz)_{kl}a_l.
  \label{eq:g-gaussian-cond-precision}
\end{equation}
For an interior frame, only the two neighbouring entries of row $k$ are
nonzero, so
\begin{align}
  \E[a_k\mid a_{\setminus k}]
  &=
  -\frac{\sigma_\eta^2}{1+\alpha^2}
  \left(
    -\frac{\alpha}{\sigma_\eta^2}a_{k-1}
    -\frac{\alpha}{\sigma_\eta^2}a_{k+1}
  \right)\\
  &=
  \boxed{
  \frac{\alpha}{1+\alpha^2}
  \bigl(a_{k-1}+a_{k+1}\bigr)
  }.
  \label{eq:g-cond-mean}
\end{align}
All information supplied by the clean chain about $a_k$ is therefore
mediated by its two neighbours.

This exact locality will not survive the diffusion channel at the level of
the observed variables. Although the clean trajectory is Markov, after
adding independent Gaussian noise and marginalising the latent clean states,
the noisy precision is generally dense. Distant noisy frames can therefore
contribute to the optimal estimate, and hence to the joint score, of a given
frame. This loss of sparsity is one of the central effects of diffusion that
we study next.


\section{The noisy law}
\label{sec:g-noisy}

Each frame is corrupted independently by the Ornstein--Uhlenbeck channel of
Section~\ref{sec:ou}. Solving the OU SDE \eqref{eq:ou-sde} once
(Toolbox~\ref{tb:ou-solution}) gives the Gaussian transition
$x_k \mid a_k \sim \Nd\big(e^{-t} a_k,\, \Deltat\big)$,
$\Deltat = 1 - e^{-2t}$, which is \eqref{eq:ou-channel}. Stacking the $L$
frames,
\begin{equation}
  x = e^{-t} a + \sqrt{\Deltat}\,\xi,
  \qquad
  \xi \sim \Nd(0, I_L),
  \quad \xi \perp a ,
  \label{eq:g-channel-vec}
\end{equation}
with $\xi$ white because the channels use independent noises.


\begin{proposition}[Noised law]\label{prop:g-Sigt}
Under the Ornstein--Uhlenbeck channel,
the noisy trajectory remains Gaussian:
\begin{equation}
  P_t(x)
  =
  \Nd(x;0,\Sigt),
  \qquad
  \boxed{
  \Sigt
  =
  e^{-2t}\Sigz+\Deltat I_L
  }.
  \label{eq:g-Sigt}
\end{equation}
\end{proposition}

\begin{proof}
Both $a$ and $\xi$ are Gaussian and independent, so their linear combination
$x=e^{-t}a+\sqrt{\Deltat}\,\xi$ is again Gaussian. Its mean is
\begin{align}
  \E[x]
  &=
  e^{-t}\E[a]
  +
  \sqrt{\Deltat}\,\E[\xi]
  =0.
\end{align}
For the covariance,
\begin{align}
  \operatorname{Cov}(x)
  &=
  \operatorname{Cov}\!\left(
    e^{-t}a+\sqrt{\Deltat}\,\xi
  \right)\\
  &=
  e^{-2t}\operatorname{Cov}(a)
  +
  \Deltat\,\operatorname{Cov}(\xi)
  +
  e^{-t}\sqrt{\Deltat}\,
  \Big(
    \operatorname{Cov}(a,\xi)+\operatorname{Cov}(\xi,a)
  \Big).
\end{align}
Independence gives
$\operatorname{Cov}(a,\xi)=0$, while
$\operatorname{Cov}(a)=\Sigz$ and $\operatorname{Cov}(\xi)=I_L$. Hence
\begin{equation}
  \operatorname{Cov}(x)
  =
  e^{-2t}\Sigz+\Deltat I_L,
\end{equation}
which proves \eqref{eq:g-Sigt}.
\end{proof}

Because
\[
  (\Sigz)_{ij}=\alpha^{|i-j|},
\]
the covariance can also be read entrywise:
\begin{equation}
  (\Sigt)_{ij}
  =
  e^{-2t}\alpha^{|i-j|}
  +
  \Deltat\,\delta_{ij},
  \label{eq:g-Sigt-entrywise}
\end{equation}
where $\delta_{ij}$ is the Kronecker delta. This separates the effect of the
channel on variances and correlations.

On the diagonal, using $\Deltat=1-e^{-2t}$,
\begin{align}
  (\Sigt)_{ii}
  &=
  e^{-2t}(\Sigz)_{ii}+\Deltat\\
  &=
  e^{-2t}+1-e^{-2t}
  =1.
\end{align}
Thus every frame remains marginally $\Nd(0,1)$ at every diffusion time: the
channel is variance preserving.

For two distinct frames, $i\neq j$, the added white noise contributes
nothing because $(I_L)_{ij}=0$. Therefore
\begin{equation}
  (\Sigt)_{ij}
  =
  e^{-2t}\alpha^{|i-j|},
  \qquad i\neq j.
  \label{eq:g-Sigt-offdiag}
\end{equation}
Since the marginal variances remain one, these entries are also the
correlations:
\begin{equation}
  \operatorname{Corr}(x_i,x_j)
  =
  e^{-2t}\alpha^{|i-j|}.
  \label{eq:g-noisy-correlation}
\end{equation}
The channel therefore leaves each one-frame distribution unchanged while
exponentially attenuating the dependence between different frames.

In matrix form,
\begin{equation}
  \Sigt
  =
  \begin{pmatrix}
    1
    & e^{-2t}\alpha
    & e^{-2t}\alpha^2
    & \cdots \\
    e^{-2t}\alpha
    & 1
    & e^{-2t}\alpha
    & \ddots \\
    e^{-2t}\alpha^2
    & e^{-2t}\alpha
    & 1
    & \ddots \\
    \vdots
    & \ddots
    & \ddots
    & \ddots
  \end{pmatrix}.
  \label{eq:g-Sigt-matrix}
\end{equation}
Thus $\Sigt$ remains symmetric Toeplitz. For $\alpha\neq0$ it is dense at
every finite $t$: diffusion does not create new zeros in the covariance,
but uniformly shrinks every off-diagonal entry by $e^{-2t}$.

The two limiting cases make the interpolation explicit:
\begin{align}
  \Sigt\big|_{t=0}
  &=
  \Sigz,\\
  \lim_{t\to\infty}\Sigt
  &=
  I_L.
  \label{eq:g-Sigt-limits}
\end{align}
At $t=0$ we recover the clean correlated trajectory; as $t\to\infty$, all
cross-frame correlations vanish and the frames approach independent standard
Gaussians.

The important change, however, appears in the \emph{precision}, not directly
in the covariance. Define
\begin{equation}
  \Qt:=\Sigt^{-1}.
  \label{eq:g-Qt-def}
\end{equation}
At $t=0$,
\[
  \Qt=\Qz,
\]
which is tridiagonal by Proposition~\ref{prop:g-Q0}. For positive diffusion
time, however, adding $\Deltat I_L$ to the covariance and then inverting does
not preserve that tridiagonal structure: in the nontrivial correlated case
$\alpha\neq0$, $\Qt$ is generally dense. Hence noisy frames that were
conditionally independent at the clean level become conditionally coupled
after the latent clean trajectory has been corrupted and marginalised.

This distinction is central:
\begin{equation}
  \boxed{
  \begin{array}{c}
    \text{clean precision }\Qz
    \text{ is local and tridiagonal}
    \\[3pt]
    \Downarrow\ \text{diffusion and marginalisation}
    \\[3pt]
    \text{noisy precision }\Qt
    \text{ is generally nonlocal and dense}.
  \end{array}}
  \label{eq:g-locality-loss-summary}
\end{equation}
Section~\ref{sec:g-tridiag-loss} quantifies how this loss of locality develops
with diffusion time.


\section{The score, two ways}\label{sec:g-score}

\subsection{Linear algebra}

Since $\log P_t(x) = -\tfrac12 x^\T \Qt x + \mathrm{const}$ and
$\nabla_x\big(x^\T \Qt x\big) = 2\Qt x$ for symmetric $\Qt$,
\begin{equation}
  \boxed{\;\score(x, t) \;=\; \nabla_x \log P_t(x) \;=\; -\,\Qt\, x,
  \qquad \Qt := \Sigt^{-1}.\;}
  \label{eq:g-score-matrix}
\end{equation}
The joint score is linear in $x$ --- the signature of a Gaussian $P_t$ --- and
its whole content is $\Qt$: which frames the denoiser of frame $k$ draws on,
and how strongly, is row $k$ of $\Qt$. Already at $L = 2$, with
$r := \alpha\, e^{-2t}$, \eqref{eq:g-score-matrix} gives
$\score_0(x,t) = -(x_0 - r\,x_1)/(1-r^2)$: the denoiser of frame $0$ uses frame
$1$ with weight $\alpha$ at $t = 0$, decaying as $e^{-2t}$, while the per-frame
marginal score stays $-x_k$ at every $t$ and carries nothing about the
dynamics. The lifecycle of that coupling, for all $L$, is
Sections~\ref{sec:g-lifecycle}--\ref{sec:g-bulk}.


\subsection{Bayesian inference}

The score can be obtained in two equivalent ways. Since the noisy law is
Gaussian, the previous section gives it directly using linear algebra.
There is, however, a second route that is more useful beyond the Gaussian
case: infer the clean trajectory $a$ from the noisy observation $x$, and
express the score through its posterior mean.

For the Ornstein--Uhlenbeck channel,
the score--posterior identity \eqref{eq:tweedie-joint} gives, componentwise,
\begin{equation}
  \boxed{
  \score_k(x,t)
  =
  \frac{e^{-t}\E[a_k\mid x]-x_k}{\Deltat}
  }.
  \label{eq:g-tweedie}
\end{equation}
Thus computing the score is equivalent to computing the posterior mean of
the clean chain.

\begin{proposition}[Posterior of the clean chain]
\label{prop:g-posterior}
The posterior distribution of the clean trajectory is Gaussian,
\begin{equation}
  P(a\mid x)
  =
  \Nd(a;m,J^{-1}),
\end{equation}
with posterior precision
\begin{equation}
  \boxed{
  J
  =
  \Qz+\frac{e^{-2t}}{\Deltat}I_L
  }
  \label{eq:g-posterior-precision}
\end{equation}
and posterior field
\begin{equation}
  \boxed{
  h
  =
  \frac{e^{-t}}{\Deltat}x,
  \qquad
  Jm=h.
  }
  \label{eq:g-posterior}
\end{equation}
Equivalently,
\begin{equation}
  m
  =
  \E[a\mid x]
  =
  J^{-1}h.
  \label{eq:g-posterior-mean}
\end{equation}
\end{proposition}

\begin{toolbox}{Posterior precision and field by completing the square}

Bayes' rule gives
\begin{equation}
  P(a\mid x)
  \propto
  P_0(a)\,P(x\mid a).
\end{equation}
The prior is
\begin{equation}
  P_0(a)
  \propto
  \exp\!\left(
    -\frac12 a^{\T}\Qz a
  \right),
\end{equation}
while the channel likelihood is
\begin{align}
  P(x\mid a)
  &\propto
  \exp\!\left[
    -\frac{1}{2\Deltat}
    \norm{x-e^{-t}a}^2
  \right].
\end{align}
Multiplying the two densities amounts to adding their negative
logarithms:
\begin{align}
  -2\log P(a\mid x)
  &=
  a^{\T}\Qz a
  +
  \frac{1}{\Deltat}
  (x-e^{-t}a)^{\T}(x-e^{-t}a)
  +\mathrm{const}.
  \label{eq:g-post-expand-1}
\end{align}
Expanding the likelihood term,
\begin{align}
  (x-e^{-t}a)^{\T}(x-e^{-t}a)
  &=
  x^{\T}x
  -2e^{-t}x^{\T}a
  +e^{-2t}a^{\T}a.
\end{align}
The term $x^{\T}x$ is constant with respect to $a$ and can therefore be
absorbed into the normalising constant. Hence
\begin{align}
  -2\log P(a\mid x)
  &=
  a^{\T}\Qz a
  +
  \frac{e^{-2t}}{\Deltat}a^{\T}a
  -
  2\frac{e^{-t}}{\Deltat}x^{\T}a
  +\mathrm{const}
  \\
  &=
  a^{\T}
  \left(
    \Qz+\frac{e^{-2t}}{\Deltat}I_L
  \right)a
  -
  2
  \left(
    \frac{e^{-t}}{\Deltat}x
  \right)^{\T}a
  +\mathrm{const}.
  \label{eq:g-post-expand-2}
\end{align}

A Gaussian written in information form has exponent
\begin{equation}
  -\frac12
  \left(
    a^{\T}Ja-2h^{\T}a
  \right).
\end{equation}
Matching terms therefore identifies
\begin{align}
  J
  &=
  \Qz+\frac{e^{-2t}}{\Deltat}I_L,
  \label{eq:g-J-identified}\\
  h
  &=
  \frac{e^{-t}}{\Deltat}x.
  \label{eq:g-h-identified}
\end{align}
Completing the square makes the posterior mean explicit:
\begin{align}
  a^{\T}Ja-2h^{\T}a
  &=
  (a-J^{-1}h)^{\T}
  J
  (a-J^{-1}h)
  -
  h^{\T}J^{-1}h.
\end{align}
The last term is independent of $a$, so
\begin{equation}
  P(a\mid x)
  =
  \Nd\!\left(
    a;\,
    J^{-1}h,\,
    J^{-1}
  \right),
\end{equation}
and therefore
\begin{equation}
  \boxed{
  \E[a\mid x]=J^{-1}h.
  }
\end{equation}

\medskip
\noindent\textbf{What the two terms in $J$ mean.}
The posterior precision is the sum
\begin{equation}
  \underbrace{J}_{\text{posterior precision}}
  =
  \underbrace{\Qz}_{\text{prior information}}
  +
  \underbrace{
    \frac{e^{-2t}}{\Deltat}I_L
  }_{\text{information from the channel}}.
  \label{eq:g-information-addition}
\end{equation}
The prior term couples neighbouring clean frames through the tridiagonal
matrix $\Qz$. The likelihood contributes only to the diagonal because,
conditional on $a$, each observation $x_k$ depends only on the corresponding
clean variable $a_k$. Hence $J$ remains tridiagonal for every finite
diffusion time.

This should be kept distinct from
\[
  \Qt=\Sigt^{-1},
\]
which is the precision matrix of the \emph{noisy marginal} $P_t(x)$ and is
generally dense. Thus
\[
  \boxed{
  \begin{array}{ccc}
    J
    &:&
    \text{posterior precision of }a\mid x
    \quad\text{--- tridiagonal},\\[3pt]
    \Qt
    &:&
    \text{marginal precision of }x
    \quad\text{--- generally dense}.
  \end{array}}
\]
The first retains the local Markov graph because the clean variables are
still present; the second results after those variables have been
marginalised out.
\hfill$\square$
\end{toolbox}

We can now insert the posterior mean into the score identity. Since
\begin{equation}
  m
  =
  J^{-1}h
  =
  \frac{e^{-t}}{\Deltat}J^{-1}x,
\end{equation}
equation~\eqref{eq:g-tweedie} gives
\begin{align}
  \score(x,t)
  &=
  \frac{e^{-t}m-x}{\Deltat}\\
  &=
  \left(
    \frac{e^{-2t}}{\Deltat^2}J^{-1}
    -
    \frac{1}{\Deltat}I_L
  \right)x.
  \label{eq:g-score-posterior-form}
\end{align}
This expression must equal the direct Gaussian score $-\Qt x$. The equality
can be seen without performing another matrix inversion. From
\begin{align}
  J
  &=
  \Sigz^{-1}
  +
  \frac{e^{-2t}}{\Deltat}I_L\\
  &=
  \frac{1}{\Deltat}
  \Sigz^{-1}
  \left(
    \Deltat I_L+e^{-2t}\Sigz
  \right)\\
  &=
  \frac{1}{\Deltat}
  \Sigz^{-1}\Sigt,
\end{align}
we obtain
\begin{equation}
  J^{-1}
  =
  \Deltat\,\Sigt^{-1}\Sigz
  =
  \Deltat\,\Qt\Sigz.
  \label{eq:g-Jinv-relation}
\end{equation}
Therefore
\begin{align}
  \frac{e^{-2t}}{\Deltat^2}J^{-1}
  -
  \frac{1}{\Deltat}I_L
  &=
  \frac{1}{\Deltat}
  \left(
    e^{-2t}\Qt\Sigz-I_L
  \right)\\
  &=
  \frac{1}{\Deltat}
  \left[
    \Qt
    \bigl(
      e^{-2t}\Sigz
    \bigr)
    -I_L
  \right]\\
  &=
  \frac{1}{\Deltat}
  \left[
    \Qt
    \bigl(
      \Sigt-\Deltat I_L
    \bigr)
    -I_L
  \right]\\
  &=
  -\Qt.
\end{align}
Hence the two routes agree exactly:
\begin{equation}
  \boxed{
  \frac{e^{-t}\E[a\mid x]-x}{\Deltat}
  =
  -\Sigt^{-1}x.
  }
  \label{eq:g-two-score-routes}
\end{equation}

The direct route computes the score from the noisy marginal covariance;
the Bayesian route computes it by first denoising the latent clean chain.
Their equality is exact, but the second formulation is the one that will
continue to make sense when the clean model is no longer jointly Gaussian.


\section{The spectral form}
\label{sec:g-lifecycle}

Between the tridiagonal $\Qz$ at $t = 0$ and the identity at $t = \infty$ lies
the whole lifecycle of the coupling structure. One tool resolves it: the
spectral theorem, applied to a family of matrices that is affine in $\Sigz$.

\subsection{The spectral theorem and its affine-family corollary}

\begin{theorem}[Spectral theorem, real symmetric case]\label{thm:g-spectral}
Every real symmetric $M \in \R^{n\times n}$ has an orthonormal eigenbasis:
there is an orthogonal $U$ ($U^\T U = I$) and reals $\omega_1,\dots,\omega_n$
with
\begin{equation}
  M = U \operatorname{diag}(\omega_i)\, U^\T
  = \sum_{i} \omega_i\, u_i u_i^\T ,
  \label{eq:g-spectral}
\end{equation}
the $u_i$ the columns of $U$.
\end{theorem}

\begin{corollary}[Affine family]\label{cor:g-affine}
For scalars $c, d$ and $M = U\operatorname{diag}(\omega_i)U^\T$,
\begin{equation}
  c\,M + d\,I
  = U \operatorname{diag}\big(c\,\omega_i + d\big)\, U^\T .
  \label{eq:g-affine}
\end{equation}
The eigenvectors are untouched; each eigenvalue moves along the same affine map
$\omega \mapsto c\omega + d$. If every $c\omega_i + d \neq 0$ then
$(cM + dI)^{-1} = U\operatorname{diag}\big((c\omega_i+d)^{-1}\big)U^\T$.
\end{corollary}

\emph{Proof.} $I = U U^\T = U\operatorname{diag}(1)\,U^\T$ because $U$ is
orthogonal, so $I$ is diagonal in \emph{every} orthonormal basis. Hence
$cM + dI = U\operatorname{diag}(c\omega_i)U^\T + U\operatorname{diag}(d)U^\T
 = U\operatorname{diag}(c\omega_i + d)U^\T$, and the inverse statement is the
diagonal case applied to $\operatorname{diag}(c\omega_i + d)$. \hfill$\square$

\subsection{The forward map in the eigenbasis}

The forward map \eqref{eq:g-Sigt}, $\Sigt = e^{-2t}\Sigz + \Deltat I$, is
exactly of the form \eqref{eq:g-affine} with $c = e^{-2t}$, $d = \Deltat$.

\begin{theorem}[Spectral form]\label{thm:g-three-forms}
Let $\Sigz = U \operatorname{diag}(\omega_i) U^\T$. Then for every $t \geq 0$,
\begin{equation}
  \boxed{\;
  \Sigt = U \operatorname{diag}\big( \lambda_i(t) \big) U^\T,
  \qquad
  \Qt = U \operatorname{diag}\!\Big( \frac{1}{\lambda_i(t)} \Big) U^\T,
  \qquad
  \lambda_i(t) = e^{-2t} \omega_i + \Deltat .
  \;}
  \label{eq:g-Qt-spec}
\end{equation}
\end{theorem}

\emph{Proof.} The first equality is \eqref{eq:g-affine}. Since $\Sigz \succ 0$
gives $\omega_i > 0$, every $\lambda_i(t) = e^{-2t}\omega_i + \Deltat \geq
\Deltat > 0$ for $t > 0$ (and $\lambda_i(0) = \omega_i > 0$), so the inverse
clause of Corollary~\ref{cor:g-affine} applies. \hfill$\square$

\subsection{What the spectral form means}

\emph{Diffusion never rotates the problem.} The eigenvectors of $\Sigt$ and
$\Qt$ are those of $\Sigz$, frozen for all $t$; only the eigenvalues move,
each along the scalar path $\omega_i \mapsto e^{-2t}\omega_i + \Deltat$. As
$t \to \infty$ every $\lambda_i(t) \to 1$: modes with $\omega_i > 1$ descend to
$1$, modes with $\omega_i < 1$ climb to $1$. In eigen-coordinates the noisy
covariance is an ellipsoid whose axes are fixed in direction and whose lengths
all relax to the unit sphere --- ... anisotropic at $t = 0$, isotropic at $t = \infty$
(Figure~\ref{fig:g-eigen-relax}).

\begin{figure}[t]\centering
\includegraphics[width=\linewidth]{fig_eigen_relax.pdf}
\caption[Eigen-axes fixed, eigenvalues relax to one]{The spectral form
\eqref{eq:g-Qt-spec} as geometry, for the AR(1) chain ($\alpha = 0.7$,
$L = 24$). \emph{Left:} the covariance ellipse
$\{x : x^\T \Sigt^{-1} x = 1\}$ in the eigenbasis of $\Sigz$, at four diffusion
times and in the limit. Its axes point along the fixed eigenvectors
$u_1, u_2$ at every $t$; only the semi-axes $\sqrt{\lambda_i(t)}$ change,
contracting the long axis and stretching the short one until, as
$t \to \infty$, the ellipse is the unit circle. \emph{Right:} the eigenvalue
flow $\lambda_i(t) = e^{-2t}\omega_i + \Deltat$ for all $L$ modes (grey), with
$\omega_{\max}$, $\omega_{\min}$ and a mode $\omega_i \approx 1$ highlighted.
Every $\lambda_i(t) \to 1$ --- modes above $1$ descend, modes below $1$ climb.
Diffusion never rotates the problem; it only drives the spectrum to isotropy.}
\label{fig:g-eigen-relax}
\end{figure}

\subsection{Decoupling in the eigenbasis}

Write $\tilde x = U^\T x$ for the noisy trajectory in the eigenbasis of
$\Sigz$. Because $U$ is orthogonal this is a rigid rotation --- no stretching,
no loss of information --- and it turns the joint problem into scalars.

\begin{theorem}[Decoupling]\label{thm:g-decouple}
In the coordinates $\tilde x = U^\T x$, with $\lambda_i(t)$ as in
\eqref{eq:g-Qt-spec}:
\begin{enumerate}
\item the noisy law factorises,
  $P_t(\tilde x) = \prod_{i} \Nd\big(\tilde x_i;\, 0,\, \lambda_i(t)\big)$;
\item the score splits into independent scalar components,
  $\tilde\score_i(\tilde x, t) = -\,\tilde x_i / \lambda_i(t)$, each depending
  on its own coordinate only;
\item the reverse-time denoising dynamics decouple into $L$ independent
  scalar linear Gaussian diffusions, one per mode --- OU-type but
  time-inhomogeneous, since the reverse drift coefficient
  depends on $t$ through $\lambda_i(t)$.
\end{enumerate}
\end{theorem}

\emph{Proof.} (i) $\operatorname{Cov}(\tilde x) = U^\T \Sigt U
= \operatorname{diag}(\lambda_i(t))$ by \eqref{eq:g-Qt-spec}, and an
uncorrelated Gaussian is independent. (ii)
$\tilde\score = U^\T \score(U\tilde x, t) = -\,U^\T \Qt U\, \tilde x
= -\operatorname{diag}(1/\lambda_i(t))\,\tilde x$. (iii) The reverse dynamics of
Section~\ref{sec:tweedie} are driven by the score and by isotropic noise;
under $U^\T$ the score is diagonal by (ii) and an orthogonal image of white
noise is again white, so the vector equation separates coordinate by
coordinate. \hfill$\square$

The eigenbasis is where the Gaussian chain is trivial: $L$ scalars, each a
one-dimensional denoising problem with its own signal-to-noise ratio
$e^{-2t}\omega_i/\Deltat$. Everything non-trivial in the frame basis --- the
dense $\Qt$, the band that fills and empties --- is this diagonal object seen
through the fixed rotation $U$.


\subsection{Two computational corollaries}

The covariance formula
\[
  \Sigt
  =
  e^{-2t}\Sigz+\Deltat I_L
\]
admits two useful rearrangements. They contain the same information, but
emphasise different aspects of the diffusion: the first makes the
signal-to-noise interpolation explicit, while the second explains how
non-local entries appear in the precision matrix.

\paragraph{Signal-to-noise form.}
Factor $\Deltat$ from the covariance:
\begin{align}
  \Sigt
  &=
  \Deltat
  \left(
    I_L+\frac{e^{-2t}}{\Deltat}\Sigz
  \right).
\end{align}
Defining
\begin{equation}
  \gamma_t
  :=
  \frac{e^{-2t}}{\Deltat},
  \label{eq:g-gamma-def}
\end{equation}
and inverting gives
\begin{equation}
  \boxed{
  \Qt
  =
  \frac{1}{\Deltat}
  \bigl(I_L+\gamma_t\Sigz\bigr)^{-1}.
  }
  \label{eq:g-Qt-snr}
\end{equation}
The scalar $\gamma_t$ is the ratio between the variance carried by the
attenuated signal and that injected by the channel. Since
$\Deltat=1-e^{-2t}$,
\begin{equation}
  \gamma_t
  =
  \frac{e^{-2t}}{1-e^{-2t}},
\end{equation}
so
\[
  \gamma_t\longrightarrow\infty
  \quad (t\downarrow0),
  \qquad
  \gamma_t\longrightarrow0
  \quad (t\to\infty).
\]

These limits recover the two ends of the diffusion. For small $t$,
$\gamma_t\Sigz$ dominates $I_L$, hence
\begin{align}
  \Qt
  &\approx
  \frac{1}{\Deltat}
  (\gamma_t\Sigz)^{-1}\\
  &=
  \frac{1}{\Deltat\gamma_t}\Qz\\
  &=
  e^{2t}\Qz.
  \label{eq:g-Qt-smallt}
\end{align}
Thus the noisy precision is initially close, up to an overall scale, to the
clean tridiagonal precision.

At large $t$, instead, $\gamma_t\ll1$, so
\begin{align}
  \Qt
  &\approx
  \frac{1}{\Deltat}I_L.
  \label{eq:g-Qt-larget}
\end{align}
Since $\Deltat\to1$, this tends to $I_L$. The precision is therefore local at
both extremes: tridiagonal near $t=0$, and diagonal as $t\to\infty$. The
non-local structure appears at intermediate diffusion times.

\paragraph{Resolvent form.}
A second rearrangement is more informative about how this non-locality is
created. Factor $e^{-2t}$ instead:
\begin{align}
  \Sigt
  &=
  e^{-2t}
  \left(
    \Sigz+e^{2t}\Deltat I_L
  \right).
\end{align}
Because
\begin{equation}
  e^{2t}\Deltat
  =
  e^{2t}(1-e^{-2t})
  =
  e^{2t}-1,
\end{equation}
define
\begin{equation}
  c_t:=e^{2t}-1.
  \label{eq:g-ct-def}
\end{equation}
Then
\begin{align}
  \Sigt
  &=
  e^{-2t}
  \bigl(\Sigz+c_t I_L\bigr),
\end{align}
and therefore
\begin{align}
  \Qt
  &=
  e^{2t}
  \bigl(\Sigz+c_t I_L\bigr)^{-1}.
  \label{eq:g-res-step1}
\end{align}
Using $\Sigz=\Qz^{-1}$,
\begin{align}
  \Sigz+c_tI_L
  &=
  \Qz^{-1}+c_tI_L\\
  &=
  \Qz^{-1}
  \bigl(I_L+c_t\Qz\bigr).
\end{align}
Hence
\begin{equation}
  \boxed{
  \Qt
  =
  e^{2t}
  \Qz
  \bigl(I_L+c_t\Qz\bigr)^{-1}.
  }
  \label{eq:g-Qt-resolvent}
\end{equation}
Since $\Qz$ commutes with any polynomial, and therefore with
$(I_L+c_t\Qz)^{-1}$, the order of the two factors is immaterial.

The inverse in \eqref{eq:g-Qt-resolvent} can be expanded locally using a
\emph{Neumann series}. This is the matrix analogue of the scalar geometric
series
\begin{equation}
  \frac{1}{1+z}
  =
  1-z+z^2-z^3+\cdots,
  \qquad |z|<1.
\end{equation}
For a matrix $A$, the corresponding identity is
\begin{equation}
  (I_L+A)^{-1}
  =
  \sum_{m=0}^{\infty}(-A)^m
  =
  I_L-A+A^2-A^3+\cdots,
  \label{eq:g-neumann-general}
\end{equation}
whenever the series converges. Applying it with
$A=c_t\Qz$ gives
\begin{align}
  (I_L+c_t\Qz)^{-1}
  &=
  I_L
  -c_t\Qz
  +c_t^2\Qz^2
  -c_t^3\Qz^3
  +\cdots.
\end{align}
Substituting into \eqref{eq:g-Qt-resolvent},
\begin{align}
  \Qt
  &=
  e^{2t}
  \Qz
  \left(
    I_L
    -c_t\Qz
    +c_t^2\Qz^2
    -\cdots
  \right)\\
  &=
  \boxed{
  e^{2t}
  \sum_{m=0}^{\infty}
  (-c_t)^m\Qz^{\,m+1}
  }\\
  &=
  e^{2t}
  \left(
    \Qz
    -c_t\Qz^2
    +c_t^2\Qz^3
    -\cdots
  \right).
  \label{eq:g-Qt-res}
\end{align}

For symmetric positive-definite $\Qz$, this expansion converges whenever
\begin{equation}
  c_t\,\lambda_{\max}(\Qz)<1,
  \label{eq:g-neumann-condition}
\end{equation}
and in particular holds in a sufficiently small neighbourhood of $t=0$,
because $c_t=e^{2t}-1\sim2t$.

The reason for introducing the Neumann series is not merely computational.
It reveals how the dense precision is built from powers of the clean sparse
precision. Since $\Qz$ is tridiagonal,
\[
  \Qz
  \quad\text{couples frames at distance at most }1,
\]
while
\[
  \Qz^2
  \quad\text{can couple frames at distance at most }2,
\]
and, more generally,
\begin{equation}
  (\Qz^m)_{ij}=0
  \qquad\text{whenever}\qquad
  |i-j|>m.
  \label{eq:g-power-bandwidth}
\end{equation}
Thus the expansion
\begin{equation}
  \Qt
  =
  e^{2t}
  \left(
    \Qz
    -c_t\Qz^2
    +c_t^2\Qz^3
    -\cdots
  \right)
\end{equation}
has a direct spatial interpretation. The first term contains only
nearest-neighbour couplings; the second can create next-nearest-neighbour
couplings; the third can reach one frame farther again; and so on.

For two frames at graph distance $d=|i-j|$, the first power capable of
connecting them is $\Qz^d$. Since $\Qz^d$ first appears with coefficient
$c_t^{\,d-1}$, their precision coupling is generated, near $t=0$, only at
order
\begin{equation}
  (\Qt)_{ij}
  =
  O\!\left(c_t^{\,d-1}\right)
  =
  O\!\left(t^{\,d-1}\right),
  \qquad d\geq1.
  \label{eq:g-band-fill-order}
\end{equation}
The dense noisy precision therefore does not appear everywhere at once:
diffusion fills the bands outward from the diagonal, one graph distance at
a time. This is the band-fill law studied in the next section.


\section{The lifecycle of the precision: losing tridiagonality}
\label{sec:g-tridiag-loss}

Sparsity is a property of a matrix in a basis. In the eigenbasis the
problem is diagonal at all times; in the frame basis, the operationally
meaningful one, the zeros of $\Qz$ are a cancellation that diffusion
destroys. The destruction is quantifiable, and it answers three questions:
when tridiagonality is lost, how, and how much structure survives at
each $t$. Figure~\ref{fig:g-sparsity-lifecycle} shows the whole arc at a
glance.

\begin{figure}[t]\centering
\includegraphics[width=\linewidth]{fig_sparsity_lifecycle.pdf}
\caption[The sparsity lifecycle of $\Sigt$ and $\Qt$]{Magnitude heatmaps
(log grey scale) of the noisy covariance $|\Sigt|$ (top) and precision
$|\Qt|$ (bottom), $\alpha = 0.85$, $L = 24$, at five diffusion times with
their signal-to-noise ratios $\gamma_t = e^{-2t}/\Deltat$. At $t = 0$ the
covariance is dense and the precision is tridiagonal. As $t$ grows the
covariance contracts onto its diagonal while the precision band \emph{fills}
outward, then fades: by $\gamma_t \ll 1$ both are within $O(e^{-2t})$ of the
identity.}
\label{fig:g-sparsity-lifecycle}
\end{figure}

\paragraph{When: immediately.} For $\alpha \neq 0$, every off-diagonal
entry of $\Qt$ is non-zero at every $t > 0$: no exact zero survives any
positive diffusion time. Analyticity alone does not give this --- each entry is
analytic in $t$ by \eqref{eq:g-Qt-spec}, but an analytic function that is not
identically zero can still vanish at isolated times, so the band-fill law below
only guarantees non-zero entries \emph{near} $t = 0$. The all-$t$ statement
needs the structure of the resolvent. From \eqref{eq:g-Qt-resolvent},
\begin{equation}
  \Qt
  = \frac{e^{2t}}{c_t}\Big[(I + c_t\Qz) - I\Big](I + c_t\Qz)^{-1}
  = \frac{e^{2t}}{c_t}\Big[I - (I + c_t\Qz)^{-1}\Big],
  \label{eq:g-Qt-inverse-form}
\end{equation}
so for $i \neq j$ the entry $(\Qt)_{ij}$ is $-e^{2t}/c_t$ times an
off-diagonal entry of $B^{-1}$, $B := I + c_t\Qz$. Now $B$ is symmetric
positive definite and tridiagonal with off-diagonal entries
$-c_t\alpha/\sigma_\eta^2 \neq 0$: an \emph{irreducible} tridiagonal
matrix. The explicit inverse formula for such matrices gives, for $i \le j$,
\begin{equation}
  (B^{-1})_{ij}
  = (-1)^{i+j}\,
    \frac{\theta_{i-1}\,\phi_{j+1}}{\theta_{L-1}}
    \prod_{k=i}^{j-1} b_k ,
  \label{eq:g-tridiag-inverse}
\end{equation}
where $b_k$ are the off-diagonal entries and $\theta_m$, $\phi_m$ are the
leading and trailing principal minors of $B$ --- all strictly positive because
$B \succ 0$. Every factor is non-zero, so every entry of $B^{-1}$, and hence
every off-diagonal entry of $\Qt$, is non-zero for every $t > 0$. The formula
also gives the sign pattern: for $\alpha > 0$ all off-diagonal entries of
$\Qt$ are negative (attractive couplings at every distance), while for
$\alpha < 0$ they alternate in sign with the distance $|i-j|$. At
$\alpha = 0$, by contrast, $\Sigt = (e^{-2t} + \Deltat)I = I$, so $\Qt = I$
and every density claim is empty.

\paragraph{How: one diagonal per order in $t$.} The loss is nevertheless
graded, and the resolvent form gives the exact rate at which each diagonal
switches on --- visible in Figure~\ref{fig:g-sparsity-lifecycle} as the
precision band widening by roughly one diagonal between successive panels.

\begin{theorem}[Band-fill leading order]\label{thm:g-bandfill}
For small $t$ and frame distance $d \geq 1$,
\begin{equation}
  \boxed{\;
  (\Qt)_{i, i+d}
  \;=\; (-1)^{d-1}\, (2t)^{d-1}\, \big( \Qz^{\,d} \big)_{i, i+d}
  \;+\; O(t^{d}) .
  \;}
  \label{eq:g-bandfill}
\end{equation}
\end{theorem}

\begin{toolbox}{Band fill from the Neumann series}
In the resolvent form \eqref{eq:g-Qt-res}, substitute
$c_t = 2t + 2t^2 + O(t^3)$ and $e^{2t} = 1 + O(t)$:
\begin{equation*}
  \Qt = \big(1 + O(t)\big)
  \Big[ \Qz - c_t \Qz^2 + c_t^2 \Qz^3 - \cdots \Big].
\end{equation*}
$\Qz$ is tridiagonal, so $\Qz^{\,m}$ has bandwidth $m$: the entry at
distance $d$ receives its first contribution from the single term $m = d$,
with prefactor $(-c_t)^{d-1} = (-1)^{d-1}(2t)^{d-1}(1 + O(t))$; every
later term is $O(t^d)$. Each additional frame of coupling range costs one
power of $t$: the band fills outward, one diagonal per order.
\hfill$\square$
\end{toolbox}


\paragraph{How much: the off-tridiagonal mass.}
Structure is measured here with the Frobenius norm: for $A \in \R^{L\times L}$,
\begin{equation}
  \|A\|_F \;:=\; \Big(\sum_{i,j=0}^{L-1} A_{ij}^2\Big)^{\!1/2}
  \;=\; \sqrt{\operatorname{tr}(A^\T A)}\, ,
  \label{eq:g-frob}
\end{equation}
the Euclidean length of $A$ read as a list of its $L^2$ entries. The one
property used below is additivity over disjoint entry sets: if $A = B + C$ with
$B$ and $C$ supported on disjoint index pairs, then
$\|A\|_F^2 = \|B\|_F^2 + \|C\|_F^2$.

Split $\Qt$ by that rule along the sparsity pattern of the clean precision. Let
$\Pi$ be the projection onto the \emph{Markov band} --- the tridiagonal set
$\{(i,j) : |i-j| \le 1\}$ that $\Qz$ occupies,
\begin{equation}
  (\Pi A)_{ij} \;=\; A_{ij}\,\mathbf{1}\{\,|i-j| \le 1\,\},
  \label{eq:g-bandproj}
\end{equation}
so $\Pi\Qt$ is the nearest-neighbour
approximation of the score operator and $\Qt - \Pi\Qt$ is everything it is
forced to drop. The two are supported on disjoint entries, so \eqref{eq:g-frob}
splits $\Qt$ into three orthogonal pieces,
\begin{equation}
  \|\Qt\|_F^2
  = \underbrace{\|\operatorname{diag}\Qt\|_F^2}_{\text{diagonal}}
  + \underbrace{\|\Pi\Qt - \operatorname{diag}\Qt\|_F^2}_{\text{nearest-neighbour band}}
  + \underbrace{\|\Qt - \Pi\Qt\|_F^2}_{\text{off-tridiagonal}} ,
  \label{eq:g-frob-split}
\end{equation}
and the two coupling pieces are the curves of
Figure~\ref{fig:g-markov-survival}, each divided by the clean nearest-neighbour
mass
$\|\Pi\Qz - \operatorname{diag}\Qz\|_F = \sqrt{2(L-1)}\;|\alpha|/\sigma_\eta^2$.
The object of interest is the \emph{off-tridiagonal mass}
\begin{equation}
  M_\perp(t) \;:=\; \|\Qt - \Pi\Qt\|_F
  \;=\; \Big(\sum_{|i-j|\ge 2}(\Qt)_{ij}^2\Big)^{\!1/2},
  \label{eq:g-offmass}
\end{equation}
the norm of exactly the couplings a Markov model may not carry.

At small $t$ the band-fill law \eqref{eq:g-bandfill} gives
$(\Qt)_{i,i+2} = -\,2t\,(\Qz^{2})_{i,i+2} + O(t^2)$ and
$(\Qt)_{i,i+d} = O(t^{\,d-1})$ for $d \ge 3$, so $M_\perp(t)$ leaves zero
linearly in $t$, dominated by the $d = 2$ diagonal. At large $t$, write
$\Sigt = I + e^{-2t}(\Sigz - I)$ and expand the inverse --- the Neumann series
converges once $e^{-2t}(\omega_{\max} - 1) < 1$, i.e.\ for all large enough
$t$ ---
\begin{equation}
  \Qt = I - e^{-2t}(\Sigz - I) + e^{-4t}(\Sigz - I)^2 + O(e^{-6t}),
  \label{eq:g-larget}
\end{equation}
so $\|\Qt - I\|_F \simeq e^{-2t}\,\|\Sigz - I\|_F$: every deviation from the
identity, off-band and on-band alike, decays at the single rate $e^{-2t}$, and
the score approaches the isotropic $-x$. In between, $M_\perp(t)$ rises, peaks,
and falls (Figure~\ref{fig:g-markov-survival}): the noisy sequence is most
non-Markov at intermediate diffusion times, yet even at the peak $M_\perp$
stays below one-tenth of the clean nearest-neighbour mass. The lifecycle in one line: tridiagonal at $t = 0$; densifying
at rate $t^{\,d-1}$ per diagonal $d$; maximally coupled --- in the Frobenius
measure used here, at the displayed $(\alpha, L)$ --- at intermediate $t$;
isotropic as $t \to \infty$ at rate $e^{-2t}$.


\begin{figure}[t]\centering
\includegraphics[width=\linewidth]{fig_markov_survival.pdf}
\caption[How much Markov structure survives diffusion]{Frobenius lifecycle of
the precision's inter-frame coupling ($\alpha = 0.9$, $L = 48$); every mass is
a Frobenius norm of a set of diagonals of $\Qt$, normalised by the
nearest-neighbour coupling mass of $\Qz$. The nearest-neighbour (Markov) mass
(green) decays monotonically, $\propto e^{-2t}$ once noise dominates. The
off-tridiagonal mass (blue) rises from zero, peaks at $t^\star \approx 0.05$
--- the densest precision, where the sequence is least Markov --- then decays
with the rest of the structure. Its peak is below $10\%$ of the original
coupling: in this norm, at this $(\alpha, L)$, the Markov band carries almost
all of the score operator's inter-frame coupling mass at every diffusion
time.}
\label{fig:g-markov-survival}
\end{figure}